\documentclass[12pt]{article}
\usepackage[utf8]{inputenc}
\usepackage[spanish]{babel}
\usepackage{amsmath}
\usepackage{amsfonts}
\usepackage{amssymb}
\usepackage{geometry}
\geometry{letterpaper, margin=2.5cm}

\begin{document}

\begin{center}
    \Large \textbf{Evaluación de Examen 1: Unidad I} \\
    \large Física para Ciencias de la Salud (005-1713) \\
    \vspace{0.5cm}
    \textbf{Estudiante:} Débora Echeverria \quad \textbf{C.I.:} 32.585.229
    \vspace{0.3cm} \\
    \large \textbf{Calificación Final: 9/10 puntos}
\end{center}

\vspace{0.5cm}

\section*{Análisis de Resultados}

\subsection*{1. Ángulo plano (Conversión a radianes)}
\textbf{Procedimiento y resultado:} Plantea una regla de tres basándose en la equivalencia de $180^\circ \rightarrow \pi \text{ rad}$. Ejecuta la simplificación de la fracción $\left(\frac{75}{180}\right)$ de forma directa dividiendo entre 15, obteniendo el valor exacto de $\frac{5\pi}{12}$ rad. Procedimiento claro y sin fallos. \\
\textbf{Veredicto:} \textbf{Correcto} (2/2 pts).

\subsection*{2. Sistemas de coordenadas (Longitud del hueso)}
\textbf{Procedimiento y resultado:} Separa metódicamente el cálculo de las componentes: calcula la distancia horizontal ($x = 7-1 = 6$) y la vertical ($y = 10-2 = 8$). Luego aplica el Teorema de Pitágoras obteniendo $c = \sqrt{36+64} = 10$. El procedimiento analítico es excelente y está acompañado por una gráfica completamente precisa. \\
\textbf{Veredicto:} \textbf{Correcto} (2/2 pts).

\subsection*{3. Vectores (Fuerza resultante)}
\textbf{Procedimiento y resultado:} La estudiante demuestra mucho orden planteando la suma vectorial ($\vec{F}_R = \vec{F}_1 + \vec{F}_2$). Calcula de forma separada las componentes, restando correctamente en el eje x ($45 - 15 = 30$) y sumando en el eje y ($25 + 35 = 60$). Ensambla el vector resultante sin errores: $30\hat{i} + 60\hat{j}$ N. \\
\textbf{Veredicto:} \textbf{Correcto} (2/2 pts).

\subsection*{4. Potencias de diez (Notación científica y prefijo)}
\textbf{Procedimiento y resultado:} Transforma la cifra correctamente a notación científica ($1,25 \times 10^{-5}$), indicando claramente con una flecha el desplazamiento de la coma. Sin embargo, no completó la instrucción adicional del enunciado que consistía en identificar e indicar el prefijo adecuado del Sistema Internacional (micrómetros). \\
\textbf{Veredicto:} \textbf{Incompleto / Parcialmente correcto} (1/2 pts).

\subsection*{5. Sistemas de unidades (Conversión de densidad)}
\textbf{Procedimiento y resultado:} Utiliza los factores de conversión en cadena de forma impecable, colocando los factores fraccionarios exactos tanto para la masa como para el volumen. Simplifica el planteamiento tachando las unidades adecuadamente para llegar a la unidad final deseada. El resultado que expresa ($1030 \text{ kg/m}^3$) es exacto y correcto, a pesar de un pequeño apunte residual confuso en la parte inferior de la hoja que no afectó el desarrollo principal. \\
\textbf{Veredicto:} \textbf{Correcto} (2/2 pts).

\vspace{0.5cm}
\section*{Comentarios Generales y Sugerencias}
¡Muy buen examen! La estudiante muestra un nivel analítico destacable, con desarrollos paso a paso muy limpios que facilitan el seguimiento de su lógica.
\begin{itemize}
    \item Se recomienda enfocar un poco más la atención a la lectura detallada de los enunciados teóricos para asegurar que las respuestas satisfagan todos los requerimientos solicitados, como ocurrió con el prefijo en la Pregunta 4.
\end{itemize}

\end{document}